\documentclass[12pt]{article}
\usepackage{multicol}
\usepackage{titlesec}
% passes different options than styledArticle

% Design
\usepackage{xcolor,amsmath}
\usepackage[labelfont=bf]{caption}
\usepackage[margin=0.6in]{geometry}
\usepackage[skip=4pt, indent=0pt]{parskip}
\usepackage[normalem]{ulem}
\renewcommand\labelitemi{$\bullet$}
\usepackage{graphicx}
\graphicspath{ {./} }

\usepackage[colorlinks]{hyperref}
\definecolor{mgreen}{RGB}{25, 160, 50}
\definecolor{mblue}{RGB}{30, 60, 200}
\usepackage{hyperref}
\hypersetup{
	colorlinks=true,
	citecolor=mgreen,
	linkcolor=black,
	urlcolor=mblue,
	pdftitle={Document by Shahar Perets},
	%	pdfpagemode=FullScreen,
}
\usepackage{yfonts}
\def\gothstart#1{\noindent\smash{\lower3ex\hbox{\llap{\Huge\gothfamily#1}}}
	\parshape=3 2.0em \dimexpr\hsize-2.4em 2.4em \dimexpr\hsize-2.5em 0pt \hsize}


% Titles are custom and differ from the title setup in megasum.cls
\newcommand\mc[1]  {$^{\text{\cite{#1}}}$}
\newcommand\aftertitle{\ \!\!\!\!\!\!\!\!\!\!}

\newcommand \npage {\vfil {\hfil \textbf{\footnotesize {Continue On The Next Page}}} \hfil \vfil \pagebreak}

\newcommand       {\rn}[1]{
	\textup{\uppercase\expandafter{\romannumeral#1}}
}


\author{Shahar Perets}
\title{A Compression Between The Studying Material in History as a way to Deduce {\large(the)} State Intentions}
\begin{document}
	\maketitle
	
	\hfil\textit{\small Made as a part of the English Research Project for Module E}
	
%	\addcontentsline{toc}{section}{Bibliography}
	\tableofcontents
	
	% Titles -- after TOC to not mess up with its styling
	\titleformat{\section}       [runin]
		{\large}
		{\Large ({\bfseries\thesection})\ }
		{0em}
		{}
		[\dotfill]
	\titleformat{\subsection}    [block]
		{\normalsize\itshape}
		{\normalfont\large({\bfseries{\thesubsection}}) \,\,$\sim$\,\,}
		{0em}
		{}
	\titleformat{\subsubsection} [block]
		{\normalsize\bfseries}
		{\normalfont\large\bfseries{(\thesubsubsection)}}
		{0.5em}
		{}
	
	\npage
	
	\section[Introduction]{Introduction}\aftertitle
	
	\gothstart{T}his research topic is about a subject the reader is most likely greatly familiar with -- history lessons. Perhaps, the reader would have expected the existence of certain events relevant to high-school students all over the world, such as the Cold War, colonialism in the 16$^{\text{th}}$-20$^{\text{th}}$ centuries, the First and Second World Wars (written as WWI and WWII respectively), Maoism, Nazism, and Communism. However, this is not the case; in different states and countries, there is a variety of topics being taught. 
	
	The obvious question that arises as a result of this fact is, why? If some events are so important to mankind, how is it that some states are completely missing them? 
	
	Our research focuses on three main questions: 
	\begin{enumerate}
		\item What is being taught in different countries as part of their history lessons? 
		\item What does it teach about the state's intentions? 
		\item How it impacts us, in Israel? 
	\end{enumerate}
	
	This whole research project started when I asked my history teacher, Sarit, why Israel chose to ignore International and even Jewish history, if it doesn't fit the government's intentions (more information in section (3)). According to her, each nation (In Israel's case, Zionism) tries to keep its intentions when teaching history, and that it is done internationally -- not just in Israel. 
	
	{\scriptsize Note: For the sake of sounding formal, I'll use ``we'' and ``us'' even though this project was made by just one person. }
	
	\npage
	\section{A Review of the National History Curriculum in Different Countries}
	\subsection{Appendix}
	\gothstart{I}n order to answer the first research question -- ``what is being taught in different countries as part of their history lessons?'', I tried to extract information from the most obvious source of information -- the official syllabus of the state. In this process, I ran into the following problems: 
	\begin{itemize}
		\item The official syllabus of the state is usually written in the native language of the place. In cases like the USA and Israel, it wasn't a problem, but for France and Germany (a subsection that, in the end, I dismissed), I've had to rely on a translation. I mostly used Google Translate. 
		\item Sometimes the syllabus is not homogeneous across the state. Due to those reasons, I gave up on adding a document about Germany. 
	\end{itemize}
	
	Note that in Europe there are two ``routes'' -- the first one is a technical collage, and the second is a general collage. The so-called ``collage'' does not convey the American meaning of the word, but rather means a school for 7-12$^{\text{th}}$ grades. 
	
	The following section, for obvious reasons, is by far the longest. It's divided into 5 main subsections: (1) this Intro, (2) Israel, (3) France, (4) New York State, (5) England. 
	
	\subsection{Israel}
	Probably the simplest to explain out of all of the programs. It's divided up into 3 categories\mc{israelEDUnew}: 
	\begin{multicols}{2}
		\begin{enumerate}
			\item WWII \& Holocaust
			\item Nationality \& Zionism
			\item Building a Country [Israel]
		\end{enumerate}
	\end{multicols}
	Just by these names, it's pretty clear how chapters 1 and 2 are structured: general history (General Nationality, World War II) and alongside that, the Jewish part of that more general event/phenomenon (Zionism/Holocaust). Chapter 3 is different as it contains no references to global events. Events after the 50$^{\text{th}}$ outside of Israel are missing too. It's hard to understand the current world order without a proper knowledge of the late 20$^{\text{th}}$ century. 
	
	Even though it may seem like the program is mostly focused on historical events related primarily to Judaism, it is a big advancement since earlier versions of the program. 10 years ago, events in WWII that are not related to the Holocaust, were not being learned. The same applies for general nationality\mc{israelEDUold}. 
	
	The problem is the misleading picture that is painted when viewing events like the Holocaust and even Jewish migration through the lens of just Zionism. For example, the Bund, the most popular party and Youth movements, is not mentioned in the books even once (!)\mc{israelDeniesBund}. Jewish Immigration as a result of non-Zionist intentions is not being talked about, etc. 
	
	
	\subsection{France}
	Perhaps the best example of a curriculum that tries to convey a feeling of being wide and comprehensive, but when you dig in, you find a clear bias towards a certain attitude. First of all, we (by ``we'' I'm referring to me only) need to understand two things about the education system in France as a whole: 
	\begin{itemize}
		\item The Final exam in History technically is an exam in both History and Geography. In practice, it has little to do with Geography. 
		\item France uses the weird European high-school system. There are differences between the technical and general schools, that I'll not cover. 
	\end{itemize}
	
	In short, the study material consists mostly of: 
	\begin{multicols}{2}
		\begin{enumerate}
			\item The French Revolution
			\item The First and the Second World War
			\item France and the world since 1945
		\end{enumerate}
	\end{multicols}
	I would want to emphasize how the French History Program views the French Revolution\mc{franceEDU}: 
	\begin{itemize}
		\item Public execution of people during the Napoleonic Wars, forcing Roman law onto Non-French citizens, and roughly a million civilians starving and dying, became ``legal imperialism'' in the History Programs. 
		\item Let's look at The Reign of Terror, a period between 1934-1944, during which approximately [source needed] 16 thousand people were murdered directly on the Guillotine (and 25K indirectly), sometimes because they were enemies of the Reign, and sometimes just to convey a feeling of fear and terror. 
		
		The general schools Program uses the word ``terror'' just 4 times in the document. The Reign of Terror gets a full 3 sentences' mention in the whole 14-page-long document. The Reign of Terror is not a part of the material for technical schools. 
		
		\item In the beginning of the revolution, hundreds of (many times innocent) aristocrats (or people that were chosen to be called aristocrats) were murdered every day. In the Program, it translates to ``the period was intense''. 
		\item The emphasis is that the Program is vastly on the national aspect of it. The word ``Nation'' and its derivatives are written 63 and 34 times in the general and technical programs, respectively. 
	\end{itemize}
	With that said, the Program rightfully criticizes the actions of France in regard to colonialism. So, while the Program has no problem criticizing France as a country (e.g., in the case of colonialism), it fails to give a merely objective view on anything related to France as a nation. This is unlike other nations like the UK. 
	
	I would want to remark the fact that the Program for General school covers events that happened in the world after 1945, such as the South African Genocide and the Cold War. With that said, a substantial focus of this chapter is on changes in France itself. 
	
	\subsection{New-York City}
	There is no central history program in the US as a whole -- education is largely dependent on the state. On a personal note, it was relatively easy to write this subsection, since (1) it's written in English, and not French (2) it's just one long, clearly sectioned document, not 38 different documents that half of them cover the same subject, (3) there are no history lessons in New York City. 
	
	Right. No History lessons. Instead of that, there is what's called ``Social Studies''. The exam in Social Studies is divided up into 3 parts: US history, World History, and ``Participation in Government'' (which is unrelated to the topic and I'll not cover it). Since the history part is jammed up with Geography, Economics, mostly ``Social Study Practices'', which is another way to say that they're being taught \textit{how to learn history}, rather than \textit{learning history}, it's clear that history will be much less significant in comparison to an actual history test. 
	
	One thing we can say in favor of the Social Studies framework, is that it's extremely broad. Every major event -- from the agricultural revolution, the Ottoman Empire, the Industrial Revolution, imperialism, WWI\&II, the Cold War, the collapse of the USSR, etc. Nothing is covered to a great extent\mc{NYScrap1}: 
	\begin{itemize}
		\item There's no mention of antisemitism in WWII (but one mention of the Holocaust). 
		\item Students will learn about perestroika and the Cold War, without learning what Communism is. 
		\item Students will learn about WWII without learning what Nazism is. Nazism, Nazis, the Nazi Party, etc., are not mentioned even one time in the whole subsection dedicated to WWII. 
		\item Some straight-up lies contained in the framework: 
		\begin{itemize}
			\item ``In response to World War II and the Holocaust, the United States played a major role in
			efforts to prevent such human suffering in the future'' (the official Program) -- a contradiction to The Holocaust in Hungary. 
			\item Students will investigate Zionism [...] and Arab Nationalism -- translated in reality into phrasing in the textbooks like ``Zionism [is an] Extreme of extreme nationalism belief that Jews need a homeland in Palestine''\mc{NYScrap2}, in contradiction to the definition of Zionism 
			\item China is communist (this is just formally true) -- ``[...] that resulted in [...]  \\-run People’s Republic of China'' (page 25). 
		\end{itemize}
	\end{itemize}
	
	\textbf{Important highlight: }even though the education is state-dependent, a general framework and skills required from social studies in general are suggested by an organization named NCSS (National Council for the Social Studies). It does not appear to publish anything related to history specifically, but rather, to describe general ``themes'' like \textit{``technology''}, \textit{``individuals and groups''}, \textit{``time and change''}, and other non-specific things that do not relate to history teaching in any way\mc{NCSScrap1}. Supposedly more than 50 states\mc{NCSScrap2}. Hence, by analyzing how terribly implemented those unified frameworks are in NYS, perhaps we can infer how ``Social Studies'' looks in the rest of the United States. 
	
	\subsection{England}
	The Education system in England is weird too. The age group 6-16 is separated into 4 stages: Key-Stage 1 to 4 (usually KS1-2 referred to as \textit{primary education} and KS3-4 as \textit{secondary education}). At the end of KS4, the students take a \textit{GCSE test}, and finish their duties as students. Some students may decide to continue into Further Education (FE, also referred to as \textit{sixth form}), taking specialized tests (called \textit{A-levels}\mc{tfIsALevel}) in usually 3-4 subjects, and optionally some more exams in different evaluating systems (e.g., International Baccalaureate, T-levels, and more). FE usually determines the profession or the student. FE targets students in the age group of 16-19. 
	
	Even though an A-test in history exists, some students don't make it into FE at all and choose to continue to other education programs, and those who do continue with FE -- usually do not choose history. Hence, we'll focus on the history GCSE curriculum, taught and learned by all of the students across the UK. History is learned in KS1-3; hence, people after the age of 14 (class 9), will not be required to learn history\mc{UKnational}. 
	
	Boldly, the curriculum of the GCSE requires ``British history must form a minimum of 40\% of the assessed content over the full course.''\mc{UKGeneral}. The actual share of non-British history, according to the history program for KS3, is much less than 60\%\mc{UKKS3}. In practice, the vast majority of learned material is related to England. British colonialism is not mandatory for schools to teach, and when it is being taught, it's as a short chapter. 
	
	
	\npage 
	
	\section{Deduction and Comparison}\aftertitle
	
	\gothstart In this section, I'll attempt to answer the second research question -- ``What does it [the history curriculum] teach about the state's intentions?'', using the information I collected in the first section. Considering all of the information has already been gathered, these last two sections will not cite many sources. 
	
	In all of the countries above, the history of the state, \textbf{and especially the nationality} of the state, has played a critical role in the curriculum. Moreover, in some places, like New York, courses in patriotism are required by law\mc{section801}. 
	Moreover: competing ideologies have left behind, for example:
	\begin{itemize}
		\item In the case of New York, communism is not taught, while communist expansion does. 
		\item In the case of Israel, the Bund and similar non-Zionist organizations are dismissed, even if they played a more significant part in history than their Zionist counterparts. 
		\item France fails to criticize actions done in the name of Nationality. 
		\item In the UK, the students will learn mostly about British history, with the lack of more ``negative'' events like the British Colonization in South Africa. 
	\end{itemize}
	
	This shows how the intentions of the state itself are expressed in what the students are taught. With that said, all programs aim to cover a wide variety of topics that are unrelated to the state itself. 
	
	%	In comparison to Israel, we found that most countries have a bit more diverse curriculum in relation to general history. Mainly, the Cold War is missing from the syllabus. Nevertheless, Israel has progressed towards a more general program in the last 7 years. 
	% well after the UK, that's not true at all
	
	One last substantial note, is that none of the programs even mentions the history of eastern society like China and Japan (apart from a small part in New York's program). The same applies to third-world countries. It makes sense considering the disconnection between the Eastern and the Western societies, but \textit{there is a possibility that it's a consequence of geopolitical tensions} (a counterpoint for that argument is the lack of information about non-aligned second-world countries, or even Arab countries that lean towards the West like the UAE, Saudi Arabia, etc.). Regardless of the reason, the nonexistence of any information about the origin and the culture of the vast majority of the people in the world, who live outside of Europe, America and Russia -- is a major problem. 
	
	We conclude that on one hand, across the studied states, a unified effort to bring the students to a place where they have a general knowledge about how most things in the world came into existence, while still leaning heavily towards teaching their own agenda. 
	
	\npage
	\section{How It Impacts Us?}\aftertitle
	
	\gothstart Considering the information in the previous section, the third research question naturally arises -- how does all of this, impact \textit{us}, in Israel? 
	
	One attempt to answer this question, may be the reason this whole project has started in the first place -- to address the lack of General History in Israel, my history teacher justified that (in part) using the argument that this is the situation across many other countries. Turns out, she was right, but I still dismiss that argument as a kind of ``Ad Hominem'' that fails to attack the substance of the argument itself. 
	
	I suggest other answers to this question. 
	\begin{itemize}
		\item \textbf{A way to understand the other}. When talking and arguing with people who are not from Israel, it's important to bear in mind what the other side might know or not know, and vice versa. The majority of people don't know much more about history than what they're being taught in school; what is assumed as forgiven to us, could not be for the other, and being aware of the curriculum will help us add the required context. For example, it's useless to argue about Nazism, if the person whom you are talking with has never heard a formal definition of the Nazi ideology, and the only thing they know about the holocaust is that Hitler killed 6 million Jewish people (as is the case in New York). 
		\item \textbf{A way to be a better student}. Every book, test, and paper, especially in not-so-well-defined subjects like history, has the bias of its writers. And in the case of national education, all of that bias is governed by the state, which has its own intentions and biases. As students, by knowing what is expected of us, we can get a more informed and educated understanding of the history. 
	\end{itemize}
	
	I think those two answers fit the last research questions better. 
	
	\npage
	\section{Conclusion}\aftertitle
	
	\gothstart We conclude that on one hand, there is a unified effort to bring the students to a place where they have a general knowledge about how most things in the world came into existence, while still leaning heavily towards teaching their own agenda. 
	
	A learner, in any country, has to always be aware and critical of what is being learned, as history can be twisted and viewed from different angles. Perhaps more importantly, it is to be conscious of what is \textit{not} learned, and why. As my history teacher once said -- every history book has an agenda. The way major events like The Holocaust, The Cold War, and The French Revolution are being learned across the globe is different! One will always try push their agenda towards you. 
	
	We learned that nationality plays a big part in \textit{what} and \textit{how} things are learned. France wants you to be proud you're French, Israel would like you to be proud you're Israeli, and so on. Those intentions are present in the history curriculum of each nation. 
	
	Another notable remark was the lack of information about the age of the Renaissance, even in countries (like France and the UK) that do teach about the Middle Ages. The transition from the Middle Ages to the modern period, notably how the countries and empires of the 16$^{\text{th}}$ and 17$^{\text{th}}$ centuries have come into existence. I can only speculate on the reasons; my guess is that the geopolitical situation was complex and fluid, and vastly different from today's. 
	
	\npage
	\section{Reflection}\aftertitle
	
	\gothstart Since I was the only member of my group, I was responsible for the entirety of it, while the most difficult aspect for me, was understanding the translation of the 38 different documents on the French Education Department website. In the French E.D., aside from the poor translation, everything was not organized well; the same information was separated into 4 documents all over the place, and when I wanted to refer back to something I read, I found it very difficult to find it in the French case. Other departments just had one or two long, clearly divided PDFs. 
	
	With that said, the project was interesting. It helped me understand why and how people in different countries grow up with specific opinions, and emphasized the importance of being critical of anything. 
	
	I don't think this project helped my English improve. The most problematic for me is speaking fluently, while the written part does not involve any speaking at all, and the Podcast is pre-written. 
	
	If I were to do the project today, I don't think I would have changed anything. It didn't make me a better independent learner, nor more confident in myself, since I'm used to work on my own in University. As for AI engines, I relied heavily on Google Translate to translate articles in French and Chinese, and I used Adobe AI tools to denoise parts of the podcast. 
	
	\npage
	
	\addcontentsline{toc}{section}{References}
	\bibliographystyle{plain}
	
	\begin{thebibliography}{1}
		\bibitem{israelEDUnew}Israel's Education Department $\sim$ (2025-26) $\sim$ \href{https://meyda.education.gov.il/files/Curriculum/history_mm_10_12.pdf}{website}
		
		\bibitem{israelEDUold}Israel's Education Department $\sim$ (2016-17) $\sim$ \href{https://meyda.education.gov.il/files/Mazkirut_Pedagogit/MikzootAzmaeem/historia_tachnit_hadasha.pdf}{website}
		
		\bibitem{franceEDU}France' National Ministry Education $\sim$ Official Bulletin 22\&25 for 2019 [still used today] $\sim$ \href{https://eduscol.education.fr/1667/programmes-et-ressources-en-histoire-geographie-voie-gt}{website} 
		
		\bibitem{israelDeniesBund}Israel Education Department Books $\sim$ \href{https://meyda.education.gov.il/files/noar/70yeras2.pdf}{website}
		
		\bibitem{NYScrap1} NYS Education Department $\sim$ New York State K-12 Social Studies framework $\sim$ (2014 [still used today, updated in 2016 and in 2017]) $\sim$ \href{https://www.nysed.gov/standards-instruction/social-studies}{website}
		
		\bibitem{NCSScrap1} The NCSS $\sim$ National Curriculum Standards for Social Studies: A Framework for Teaching, Learning, and Assessment $\sim$ (2010) $\sim$ \href{https://www.socialstudies.org/standards/national-curriculum-standards-social-studies}{website}
		
		\bibitem{NCSScrap2} The NCSS $\sim$ About National Council for the Social Studies $\sim$ \href{https://www.socialstudies.org/about}{website}
		
		\bibitem{NYScrap2} The Jerusalem Post $\sim$ 10th-grade NY review handout defines Zionism as 'extreme land-grabbing nationalism' (2025) $\sim$ \href{https://www.jpost.com/diaspora/antisemitism/article-857234}{website}
		
		\bibitem{zion} Oxford Dictionary $\sim$ Definition of Zionism $\sim$ \href{https://www.oxfordlearnersdictionaries.com/definition/english/zionism}{website}
		
		\bibitem{section801} The New York State Senate $\sim$ Section 801 $\sim$ \href{https://www.nysenate.gov/legislation/laws/EDN/801}{website}
		
		\bibitem{UKnational} The UK Government $\sim$ National curriculum in England: framework for key stages 1 to 4 (2014) $\sim$ \href{https://www.gov.uk/government/publications/national-curriculum-in-england-framework-for-key-stages-1-to-4}{website}
		
		\bibitem{UKhistory} The UK Government $\sim$ GCSE (9 to 1) subject-level conditions and requirements for history (2015) $\sim$ \href{https://www.gov.uk/government/publications/gcse-9-to-1-subject-level-conditions-and-requirements-for-history}{website}
		
		\bibitem{UKGeneral} The GCSE $\sim$ SE Subject Level Conditions and Requirements for History (2015) $\sim$ \href{https://assets.publishing.service.gov.uk/media/5a74dfa9ed915d502d6cbae0/gcse-subject-level-conditions-and-requirements-for-history.pdf}{website}
		
		\bibitem{UKKS3} United Kingdom Department of Education $\sim$ History programs of study: key stage 3 (2013) $\sim$ \href{https://assets.publishing.service.gov.uk/media/5a7c66d740f0b626628abcdd/SECONDARY_national_curriculum_-_History.pdf}{website}
		
		\bibitem{tfIsALevel} The BBC $\sim$ ``School Parents -- AS and A levels'' (2010) $\sim$ \href{https://web.archive.org/web/20100728135656/http://www.bbc.co.uk/schools/parents/alevels/}{link (web archive)}
		
	\end{thebibliography}
	
	{\large\dotfill\normalsize \\ \vfil {\begin{center}
				{\textbf{\textit{Shahar Perets, 2025}} \\
					\scriptsize Compiled In {\LaTeX}\ and made using Free Software}
		\end{center}} \vfil	}
	
	
	
	
\end{document}